\chapter{Sylow theorems, the fixture}
\section{Definitions}

\begin{definition}
Let $p$ be a prime. A \vocab{$p$-group} is a group whose order is a power
of $p$.
\end{definition}

\begin{example}
The cyclic group of order $8$ is a $2$-group.
\end{example}

\section{Statement}

\begin{theorem}[Sylow]
	\label{thm:sylow}
	A finite group whose order has $p^n$ as its largest power of $p$ has a
	subgroup of order $p^n$.
\end{theorem}
\begin{proof}
	Let the group act on the subsets of size $p^n$, and count the orbits.
	\begin{lemma}
		\label{lem:orbit-count}
		Some orbit has a size that $p$ does not divide.
	\end{lemma}
	Then \Cref{lem:orbit-count} gives a stabilizer of order $p^n$.
\end{proof}

\begin{ques}
Why does \Cref{thm:sylow} need the group to be finite?
\end{ques}

\begin{remark}
A remark is not an item.
\end{remark}

\section{\problemhead}

\begin{dproblem}
Show that a group of order $15$ is cyclic, with \Cref{thm:sylow}.
\end{dproblem}
